\begin{proposition}[A singleton support retains logarithmically many corrections]
\label{prop:aesp-cd-singleton-corrections}
Let $G$ consist of one edge, let the seed be one endpoint, and choose
\begin{equation}
 0<q\leq q_0:=\frac{\sqrt{41}-5}{8},
 \qquad
 \alpha:=\frac{q^2}{1+q^2},
 \qquad
 \frac{1-\alpha}{2}\leq\rho<1.
 \label{eq:aesp-cd-singleton-parameters}
\end{equation}
For the exact safeguarded recurrence
\cref{eq:aesp-cd-safe-outer}, the RPPR optimum and every nonzero iterate are
supported only on the seed.  Let $e_t:=x_1^\star-x_{t,1}$ denote the scalar
error on that coordinate and define
\begin{equation}
 \vartheta:=\frac{2(1-q^2)}{3},
 \qquad
 u:=\frac{(1-q)(1+4q)}{3(1+q)},
 \qquad
 \xi:=\vartheta u.
 \label{eq:aesp-cd-singleton-ratios}
\end{equation}
Then, for every $k\geq0$,
\begin{equation}
 e_{2k}=\xi^ke_0,
 \qquad
 e_{2k+1}=\vartheta\xi^ke_0.
 \label{eq:aesp-cd-singleton-errors}
\end{equation}
There is no correction at an odd stage, whereas every even stage $2k\geq2$
is a full correction:
\begin{equation}
 \begin{aligned}
 r_{2k+1}&=0,&&k\geq0,\\
 r_{2k}&=\beta_{\mathrm A}d_{2k},&&k\geq1.
 \end{aligned}
 \label{eq:aesp-cd-singleton-schedule}
\end{equation}
Thus the raw count through $T$ stages is
\begin{equation}
 N_T=\left\lfloor\frac{T-1}{2}\right\rfloor.
 \label{eq:aesp-cd-singleton-count}
\end{equation}
Nevertheless, every one of these corrections satisfies
$\mathfrak D_{2k}<0$ and $\gamma_{2k}<1$, and hence
\begin{equation}
 \mathcal I_T=0
 \qquad\text{for every }T.
 \label{eq:aesp-cd-singleton-zero-inflation}
\end{equation}

For fixed admissible $\rho$, reaching scalar error $\eta$ takes
$\Theta(\log(e_0/\eta))$ stages and incurs the same order of nonzero
correction calls.  Equivalently, since
$F_\rho(x_t)-F_\rho(x^\star)=((1+\alpha)/4)e_t^2$, reaching objective gap
$\epsobj$ incurs $\Theta(\log(1/\epsobj))$ corrections.  The seed entry and
the only edge are discovered once.  Each subsequent correction/rekey call
has degree work one, so this is cumulative correction work rather than
repeated support-discovery work.  It is a statement about the safeguarded
recurrence, not a lower bound for a different algorithm or a cached response
primitive.
\end{proposition}

\begin{proof}
Let the seed be vertex $1$.  On this graph
\begin{equation*}
 Q=\begin{pmatrix}
 (1+\alpha)/2&-(1-\alpha)/2\\
 -(1-\alpha)/2&(1+\alpha)/2
 \end{pmatrix},
 \qquad
 c_\rho:=b-\alpha\rho\one
 =\alpha\begin{pmatrix}1-\rho\\-\rho\end{pmatrix}.
\end{equation*}
Writing $\lambda:=(1+\alpha)/2$, the obstacle KKT conditions give
\begin{equation*}
 x^\star
 =\begin{pmatrix}\alpha(1-\rho)/\lambda\\0\end{pmatrix}.
\end{equation*}
Indeed, the inactive-coordinate inequality is exactly
$\rho\geq(1-\alpha)/2$.  The same direct KKT check, or Stieltjes comparison,
shows that every shifted minimizer centered below $x^\star$ also has zero
second coordinate.  Thus the entire recurrence is scalar after the one-time
edge exposure.

Put $\kappa:=1-2\alpha$.  For a scalar safe center $\ell$, exact shifted
minimization gives
\begin{equation*}
 p(\ell)=\frac{\lambda x^\star+\kappa\ell}{\lambda+\kappa},
 \qquad
 x^\star-p(\ell)=\vartheta(x^\star-\ell),
 \qquad
 \vartheta=\frac{\kappa}{\lambda+\kappa}
            =\frac{2(1-q^2)}3.
\end{equation*}
Also $\beta_{\mathrm A}=(1-q)/(1+q)$.  Starting immediately after a full
correction, one shifted solve therefore changes the error from $e$ to
$\vartheta e$.  At the following odd stage,
\begin{equation*}
 x^\star-y
 =\bigl((1+\beta_{\mathrm A})\vartheta
        -\beta_{\mathrm A}\bigr)e
 =ue>0.
\end{equation*}
There is no correction, and the next shifted solve changes the error to
$\vartheta ue=\xi e$.

At the following even stage, the extrapolate instead satisfies
\begin{equation*}
 y-x^\star
 =\vartheta\bigl(\beta_{\mathrm A}
            -(1+\beta_{\mathrm A})u\bigr)e,
 \qquad
 \beta_{\mathrm A}-(1+\beta_{\mathrm A})u
 =\frac{(1-q)(1-5q)}{3(1+q)^2}>0.
\end{equation*}
For this one-coordinate trial point, the retraction magnitude is
$\Delta=(\lambda/\alpha)(y-x^\star)$.  The condition that it removes the
entire extrapolation, $\Delta\geq\beta_{\mathrm A}d$, reduces after
cancelling positive factors to
\begin{equation*}
 4q^2+5q\leq1,
\end{equation*}
which is exactly $q\leq q_0$.  Hence $\ell=x$ at every even correction and
the same two-stage calculation repeats at every scale.  This proves
\cref{eq:aesp-cd-singleton-errors,eq:aesp-cd-singleton-schedule} and the
count in \cref{eq:aesp-cd-singleton-count}.

At every correction put $a_q:=(1-q)/q$ and
$h:=((1+q)/q)r=a_qd$.  Its defect is
\begin{equation*}
 \mathfrak D=h(2e-h)<0,
\end{equation*}
because overshoot gives $\beta_{\mathrm A}d>e$ and hence
$a_qd=((1+q)/q)\beta_{\mathrm A}d>2e$.  Noncorrection stages have
$\gamma_t=1$, proving \cref{eq:aesp-cd-singleton-zero-inflation}.  Finally,
$0<\xi<1$ and \cref{eq:aesp-cd-singleton-errors} give the logarithmic
accuracy statements.  The work statement follows because the seed vertex
has graph degree one and the shared model charges its adjacency scan.
\end{proof}

For a concrete rational instance, take $q=1/7$, so
$\alpha=1/50$, $\beta_{\mathrm A}=3/4$,
$\vartheta=32/49$, $u=11/28$, and $\xi=88/343$; every
$\rho\in[49/100,1)$ is admissible.  This exact example shows that the
collapse identity cannot remove the logarithm from the \emph{number} of
correction calls, even when the support is a singleton and never expands.
It also shows why $N_T$ must not replace $\mathcal I_T$ indiscriminately:
the normalized inflation correctly assigns zero cost to all of these
benign resets.

In the singleton family,
$\mathcal B_t^{\mathrm{col}}=\varnothing$ at every correction, so
\cref{lem:aesp-cd-collateral-charge} also proves
$\mathcal I_T=0$.  In contrast, the positive three-path defect below forces
$\mathcal B_3^{\mathrm{col}}\neq\varnothing$ and is charged by
$\mathfrak C_3^{\mathrm{col}}$.
